\documentclass[12pt]{article}
\usepackage[spanish]{babel}
\usepackage[utf8]{inputenc}
\usepackage[T1]{fontenc}
\usepackage{amsmath, amssymb}
\usepackage{helvet}
\usepackage{geometry}
\usepackage{enumitem}
\usepackage{needspace}
\usepackage{tikz}
\usetikzlibrary{babel}

\geometry{letterpaper, margin=2cm}
\setlength{\emergencystretch}{2em}
\renewcommand{\familydefault}{\sfdefault}
\setlength{\parindent}{0pt}
\setlength{\parskip}{6pt}
\hyphenpenalty=10000
\exhyphenpenalty=10000

\newcommand{\opcion}[2]{\item $#2$}
\newcommand{\puntografico}{0.09}
\newcommand{\puntodestacado}{0.16}
\newcommand{\rotuloxeste}{\mbox{x (este)}}
\newcommand{\rotuloynorte}{\mbox{y (norte)}}

\begin{document}

\begin{center}
{\large \textbf{Tercera Pr\'actica de Secundaria de la Prueba Nacional Estandarizada Sumativa de Matem\'atica 2026.}}
\end{center}

\noindent Nombre: \rule{8cm}{0.4pt}\hfill Sec: \rule{3cm}{0.4pt}

\vspace{6pt}
\noindent\textbf{INFORMACI\'ON GENERAL}

\vspace{4pt}
\noindent\textbf{Materiales necesarios para la prueba:}

\begin{itemize}[leftmargin=1.6em, itemsep=1pt, topsep=2pt, parsep=0pt]
  \item Un cuadernillo que contiene:
  \begin{itemize}[leftmargin=1.6em, itemsep=0pt, topsep=1pt, parsep=0pt, label=$\blacklozenge$]
    \item informaci\'on general
    \item 60 \'items de selecci\'on de respuesta
  \end{itemize}
  \item Hoja de respuestas para lectora \'optica
  \item Bol\'igrafo con tinta azul o negra
  \item Corrector l\'iquido blanco
\end{itemize}

\noindent\textbf{Material que se puede requerir para la prueba:}

\begin{itemize}[leftmargin=1.6em, itemsep=1pt, topsep=2pt, parsep=0pt]
  \item Calculadora
\end{itemize}

\vspace{4pt}
\noindent\textbf{Instrucciones:}

\begin{enumerate}[leftmargin=*, itemsep=3pt, topsep=2pt, parsep=0pt]
  \item La Prueba Nacional Estandarizada de secundaria de matem\'atica est\'a compuesta por 60 \'items. Verifique que el cuadernillo que tiene en sus manos est\'e bien compaginado y contenga los 60 \'items de selecci\'on de respuesta correspondientes al componente Matem\'aticas. En caso de encontrar alguna irregularidad, notif\'iquela inmediatamente al delegado de aula; de lo contrario, usted asume la responsabilidad sobre los problemas que se pudieran suscitar por esta causa.

  \item Cada \'item presenta tres posibilidades de respuesta: A), B), C) y D) Solamente una de ellas es la respuesta correcta.

  \item Lea cuidadosamente cada \'item y sus respectivas opciones. Puede utilizar el espacio al lado de cada \'item para realizar cualquier anotaci\'on que necesite, con el fin de hallar la respuesta.

  \item Ning\'un \'item debe aparecer sin respuesta o con m\'as de una marca en la hoja lectora \'optica.

  \item Una vez que haya revisado todas las opciones y tenga seguridad de su elecci\'on, rellene completamente el c\'irculo correspondiente, en la hoja lectora \'optica, tal como se indica en el siguiente ejemplo:

  \begin{center}
  \begin{tikzpicture}[baseline=-0.5ex]
    \draw[thick] (0,0) circle (0.38);
    \node at (0,0) {\textbf{A}};
    \fill (1.1,0) circle (0.38);
    \node[white] at (1.1,0) {\textbf{B}};
    \draw[thick] (2.2,0) circle (0.38);
    \node at (2.2,0) {\textbf{C}};
  \end{tikzpicture}
  \end{center}

  \item Si necesita rectificar la respuesta, utilice corrector l\'iquido blanco sobre el c\'irculo por corregir y rellene con bol\'igrafo de tinta negra o azul la nueva opci\'on seleccionada. Adem\'as, en el espacio de observaciones de la hoja lectora \'optica debe anotar y firmar la correcci\'on efectuada (Ejemplo: 12=A, firma). Se firma solo una vez al final de todas las correcciones.
\end{enumerate}

\vspace{8pt}
\noindent Creado por: Alberto Jim\'enez Ruiz asesor regional de Matem\'atica de la DRE-Cartago 2026.

\newpage

\begin{enumerate}[label=\textbf{\arabic*.}, itemsep=2.5cm]

\Needspace{0.98\textheight}
\item
En un parque se instalar\'a una c\'amara de seguridad K ubicada a $6$ m al este y $3$ m al sur de la fuente representada por el punto F. La fuente se considera como origen y las unidades est\'an en metros. La c\'amara tiene un alcance m\'aximo de $4$ metros a su alrededor.

\textquestiondown Cu\'al de las siguientes representaciones gr\'aficas corresponde al alcance m\'aximo de esa c\'amara?

\vspace{1\baselineskip}

\begin{center}
\setlength{\tabcolsep}{0pt}
\begin{tabular}{r@{\hspace{0.55cm}}c@{\hspace{1.2cm}}r@{\hspace{0.55cm}}c}
\raisebox{3.45cm}{\textbf{A)}} &
\begin{tikzpicture}[scale=0.38]
  \path[use as bounding box] (-2,-9) rectangle (13,5);
  \draw[<->, thick] (-1,0) -- (12,0);
  \draw[<->, thick] (0,-8) -- (0,4);
  \node[anchor=west, inner sep=1pt] at (12.05,0) {\scriptsize\rotuloxeste};
  \node[anchor=south, inner sep=1pt] at (0,4.25) {\scriptsize\rotuloynorte};
  \draw[thick] (6,-3) circle (4);
  \draw[dashed] (6,0) -- (6,-3);
  \draw[dashed] (10,0) -- (10,-3);
  \draw[dashed] (0,-3) -- (10,-3);
  \fill (0,0) circle (\puntodestacado) node[below left] {F};
  \fill (6,-3) circle (\puntodestacado) node[anchor=south east, inner sep=1pt] {K};
  \fill (10,-3) circle (\puntodestacado);
  \node[fill=white, anchor=north, inner sep=1pt] at (6,0) {$6$};
  \node[fill=white, anchor=north, inner sep=1pt] at (10,0) {$10$};
  \node[fill=white, anchor=west, inner sep=1pt] at (0,-3) {$-3$};
\end{tikzpicture}
&
\raisebox{3.45cm}{\textbf{B)}} &
\begin{tikzpicture}[scale=0.38]
  \path[use as bounding box] (-9,-2) rectangle (7,11);
  \draw[<->, thick] (-8,0) -- (6,0);
  \draw[<->, thick] (0,-1) -- (0,10);
  \node[anchor=west, inner sep=1pt] at (6.05,0) {\scriptsize\rotuloxeste};
  \node[anchor=south, inner sep=1pt] at (0,10.25) {\scriptsize\rotuloynorte};
  \draw[thick] (-6,3) circle (4);
  \draw[dashed] (-6,0) -- (-6,3);
  \draw[dashed] (-2,0) -- (-2,3);
  \draw[dashed] (0,3) -- (-6,3);
  \fill (0,0) circle (\puntodestacado) node[below right] {F};
  \fill (-6,3) circle (\puntodestacado) node[anchor=south west, inner sep=1pt] {K};
  \fill (-2,3) circle (\puntodestacado);
  \node[fill=white, anchor=south, inner sep=1pt] at (-6,0) {$\llap{-}6$};
  \node[fill=white, anchor=north, inner sep=1pt] at (-2,0) {$-2$};
  \node[fill=white, anchor=west, inner sep=1pt] at (0,3) {$3$};
\end{tikzpicture}
\\[1.25cm]
\raisebox{3.45cm}{\textbf{C)}} &
\begin{tikzpicture}[scale=0.38]
  \path[use as bounding box] (-2,-5) rectangle (13,9);
  \draw[<->, thick] (-1,0) -- (12,0);
  \draw[<->, thick] (0,-4) -- (0,8);
  \node[anchor=west, inner sep=1pt] at (12.05,0) {\scriptsize\rotuloxeste};
  \node[anchor=south, inner sep=1pt] at (0,8.25) {\scriptsize\rotuloynorte};
  \draw[thick] (6,3) circle (4);
  \draw[dashed] (6,0) -- (6,3);
  \draw[dashed] (10,0) -- (10,3);
  \draw[dashed] (0,3) -- (10,3);
  \fill (0,0) circle (\puntodestacado) node[below left] {F};
  \fill (6,3) circle (\puntodestacado) node[anchor=south east, inner sep=1pt] {K};
  \fill (10,3) circle (\puntodestacado);
  \node[fill=white, anchor=north, inner sep=1pt] at (6,0) {$6$};
  \node[fill=white, anchor=north, inner sep=1pt] at (10,0) {$10$};
  \node[fill=white, anchor=east, inner sep=1pt] at (0,3) {$3$};
\end{tikzpicture}
&
\raisebox{3.45cm}{\textbf{D)}} &
\begin{tikzpicture}[scale=0.38]
  \path[use as bounding box] (-2,-9) rectangle (15,5);
  \draw[<->, thick] (-1,0) -- (14,0);
  \draw[<->, thick] (0,-8) -- (0,4);
  \node[anchor=west, inner sep=1pt] at (14.05,0) {\scriptsize\rotuloxeste};
  \node[anchor=south, inner sep=1pt] at (0,4.25) {\scriptsize\rotuloynorte};
  \draw[thick] (6,-3) circle (7);
  \draw[dashed] (6,0) -- (6,-3);
  \draw[dashed] (13,0) -- (13,-3);
  \draw[dashed] (0,-3) -- (13,-3);
  \fill (0,0) circle (\puntodestacado) node[below left] {F};
  \fill (6,-3) circle (\puntodestacado) node[anchor=south east, inner sep=1pt] {K};
  \fill (13,-3) circle (\puntodestacado);
  \node[fill=white, anchor=north, inner sep=1pt] at (6,0) {$6$};
  \node[fill=white, anchor=north, inner sep=1pt] at (13,0) {$13$};
  \node[fill=white, anchor=west, inner sep=1pt] at (0,-3) {$-3$};
\end{tikzpicture}
\end{tabular}
\end{center}

\Needspace{0.90\textheight}
\item
En una zona de juegos se coloc\'o una fuente circular. La siguiente representaci\'on gr\'afica, cuyas unidades est\'an en metros, muestra las posiciones, P de un punto de referencia, F del centro de la fuente y B de un punto del borde de la fuente:

\begin{center}
\begin{tikzpicture}[scale=0.45]
  \path[use as bounding box] (-3,-3) rectangle (15,13);
  \draw[<->, thick] (-2,0) -- (14,0);
  \draw[<->, thick] (0,-2) -- (0,12);
  \node[anchor=west, inner sep=1pt] at (14.05,0) {\rotuloxeste};
  \node[anchor=south, inner sep=1pt] at (0,12.35) {\rotuloynorte};
  \draw[thick] (8,5) circle (5);
  \draw[dashed] (8,0) -- (8,5);
  \draw[dashed] (13,0) -- (13,5);
  \draw[dashed] (0,5) -- (13,5);
  \fill (0,0) circle (\puntografico) node[below left] {P};
  \fill (8,5) circle (\puntografico) node[anchor=south east, inner sep=1pt] {F};
  \fill (13,5) circle (\puntografico) node[anchor=south west, inner sep=1pt] {B};
  \node[anchor=north] at (8,0) {$8$};
  \node[anchor=north] at (13,0) {$13$};
  \node[left] at (0,5) {$5$};
\end{tikzpicture}
\end{center}

De acuerdo con la informaci\'on anterior, \textquestiondown cu\'al es la representaci\'on algebraica, cuyas unidades est\'an en metros, de la circunferencia que representa el borde de esa fuente?

\begin{enumerate}[label=\Alph*)]
  \opcion{A}{(x-8)^2+(y-5)^2=25}
  \opcion{B}{(x+8)^2+(y+5)^2=25}
  \opcion{C}{(x-8)^2+(y-5)^2=10}
  \opcion{D}{(x-13)^2+(y-5)^2=10}
\end{enumerate}

\Needspace{0.78\textheight}
\item
Una estaci\'on meteorol\'ogica registra datos mediante un sensor S ubicado a $4$ km al este y $2$ km al sur del punto de referencia de una zona de monitoreo. Ese punto de referencia se considera como origen y las unidades est\'an en kil\'ometros. El sensor cubre todos los puntos ubicados a $6$ km o menos de la estaci\'on.

Durante una prueba t\'ecnica se colocaron tres dispositivos. El dispositivo A se ubic\'o $10$ km al este y $2$ km al sur del origen; el dispositivo B, $4$ km al este y $5$ km al norte del origen; y el dispositivo C, $1$ km al oeste y $5$ km al sur del origen.

De acuerdo con la informaci\'on anterior, \textquestiondown en cu\'ales de esos dispositivos se pueden registrar datos mediante el sensor?

\begin{enumerate}[label=\Alph*)]
  \item Solo en A.
  \item Solo en A y en C.
  \item Solo en B y en C.
  \item En A, en B y en C.
\end{enumerate}

\Needspace{0.88\textheight}
\item
Una empresa de log\'istica modela la zona de cobertura de un lector de c\'odigos mediante una circunferencia. El lector se ubica en el punto L y un punto del borde de la zona de cobertura se ubica en R.

La siguiente representaci\'on gr\'afica, cuyas unidades est\'an en metros, muestra la circunferencia correspondiente a esa zona de cobertura:

\begin{center}
\begin{tikzpicture}[scale=0.42]
  \path[use as bounding box] (-10,-3) rectangle (8,14);
  \draw[<->, thick] (-9,0) -- (7,0);
  \draw[<->, thick] (0,-2) -- (0,13);
  \node[anchor=west, inner sep=1pt] at (7.05,0) {\rotuloxeste};
  \node[anchor=south, inner sep=1pt] at (0,13.35) {\rotuloynorte};
  \draw[thick] (-4,6) circle (5);
  \draw[dashed] (-4,0) -- (-4,6);
  \draw[dashed] (1,0) -- (1,6);
  \draw[dashed] (-4,6) -- (1,6);
  \fill (-4,6) circle (\puntografico) node[anchor=south east, inner sep=1pt] {L};
  \fill (1,6) circle (\puntografico) node[anchor=south west, inner sep=1pt] {R};
  \node[anchor=north] at (-4,0) {$\llap{-}4$};
  \node[anchor=north] at (1,0) {$1$};
  \node[left] at (0,6) {$6$};
\end{tikzpicture}
\end{center}

De acuerdo con la informaci\'on anterior, \textquestiondown cu\'al de las siguientes representaciones algebraicas corresponde a la circunferencia que modela la zona de cobertura del lector?

\begin{enumerate}[label=\Alph*)]
  \opcion{A}{(x+4)^2+(y-6)^2=25}
  \opcion{B}{(x-4)^2+(y+6)^2=25}
  \opcion{C}{(x-4)^2+(y-6)^2=25}
  \opcion{D}{(x-1)^2+(y-6)^2=25}
\end{enumerate}

\end{enumerate}

\end{document}
